% Front matter — versão PT-PT. Espelho de ../thesis/frontmatter/frontmatter.tex (mesmo conteúdo).

\frontmatter % numeração romana (i, ii, iii, ...) nas páginas pré-conteúdo

\pagestyle{plain}

%----------------------------------------------------------------------------------------
%	PÁGINA DE TÍTULO
%----------------------------------------------------------------------------------------

\maketitlepage

%----------------------------------------------------------------------------------------
%	DECLARAÇÃO DE INTEGRIDADE
%----------------------------------------------------------------------------------------

\begin{soi}

\vspace{1cm}

Declaro que realizei este trabalho académico com integridade.

Não plagiei nem apliquei qualquer forma de uso indevido de informação ou de falsificação de resultados
ao longo do processo que conduziu à sua elaboração.

Assim, o trabalho apresentado neste documento é original e da minha autoria, não tendo sido anteriormente
utilizado para qualquer outro fim.

Declaro ainda que tomei pleno conhecimento do Código de Conduta Ética do P.PORTO.

\vspace{0.5cm}

\textbf{Uso de Inteligência Artificial.} Em conformidade com as orientações do P.PORTO/ISEP, declaro que
foram utilizadas ferramentas de inteligência artificial durante o
desenvolvimento deste trabalho, para apoiar a redação e edição do texto e o desenvolvimento de software.
Dirigi o trabalho e revi e validei criticamente todos os resultados; as ideias, as decisões e o conteúdo
final são meus e da minha responsabilidade.

\vspace{1cm}

ISEP, Porto, \today % NOTA: substituir pela data de entrega; confirmar a redação exata exigida pela MEIA/ISEP.

\end{soi}

%----------------------------------------------------------------------------------------
%	DEDICATÓRIA (opcional)
%----------------------------------------------------------------------------------------
\begin{dedicatory}
% TODO (opcional): dedicatória.
\end{dedicatory}

%----------------------------------------------------------------------------------------
%	RESUMO (língua principal: Português)
%----------------------------------------------------------------------------------------

\begin{abstract}

Esta dissertação apresenta o InvestiGator, um sistema de alertas financeiros explicável (XAI-first) para
investidores de retalho no mercado acionista dos Estados Unidos (NYSE e NASDAQ). O sistema envia alertas via Telegram a partir
de dois gatilhos independentes (movimentos abruptos de mercado, detetados como anomalias estatísticas, e
novas notícias financeiras) e, para cada alerta, expõe uma explicação totalmente rastreável: o evento
detetado, o raciocínio aplicado, as fontes e os precedentes históricos. O seu núcleo é um motor de correlação
notícia--mercado que, dada uma notícia, recupera notícias históricas análogas de uma base de conhecimento e
mede o impacto que esses precedentes tiveram nos preços, à maneira de um estudo de evento e sem lookahead.
Enquanto contribuição de Engenharia de Inteligência Artificial, o trabalho integra, aplica e avalia
criticamente componentes existentes (recuperação por embeddings de frases, deteção estatística de anomalias,
triagem supervisionada e explicação transparente) num sistema funcional, explicável e reproduzível. Em dados reais, a recuperação
semântica superou claramente as baselines lexical e triviais (precision@5 cross-ticker de $0{,}51$, face a
$0{,}35$ de uma baseline lexical e $0{,}24$ do acaso) e o detetor normalizado pela volatilidade disparou de forma muito mais consistente entre
ações do que um limiar fixo. Um modelo de triagem de materialidade treinado em $79\,753$ exemplos
notícia--mercado de vários anos quase quadruplicou a precisão dos alertas dentro do orçamento diário em
dados retidos, mas nenhum modelo com texto bateu a baseline de só-volatilidade, e em produção o mesmo
gate não mostrou benefício significativo sobre a população, distribuída de outra forma, que de facto vê.
Ambos os resultados são reportados tal como são. As limitações e o trabalho futuro ficam explícitos.

\end{abstract}

\begin{abstractotherlanguage}

This dissertation presents InvestiGator, an explainable (XAI-first) financial-alert system for retail investors in
the US equity market (NYSE and NASDAQ). The system raises Telegram alerts from two independent triggers (abrupt
market movements, detected as statistical anomalies, and incoming financial news) and, for every alert,
exposes a fully traceable explanation: the detected event, the reasoning applied, the sources, and historical
precedents. Its core is a news--market correlation engine that, given a news item, retrieves analogous
historical news from a knowledge base and measures the impact those precedents had on prices, event-study
style and without lookahead. As an Artificial Intelligence Engineering contribution, the work
integrates, applies and critically evaluates existing components (sentence-embedding retrieval, statistical
anomaly detection, supervised triage and transparent explanation) in a functional, explainable and reproducible system. On real
data, semantic retrieval clearly outperformed lexical and trivial baselines (cross-ticker precision@5 of $0.51$
against $0.35$ lexical and $0.24$ random), and the volatility-normalised detector
fired far more consistently across stocks than a fixed threshold. A materiality-triage model trained on
$79{,}753$ multi-year news--market examples nearly quadrupled within-budget alert precision, yet no text
model beat the volatility-only baseline, a finding reported as it stands. Limitations and future work are
stated explicitly.

\end{abstractotherlanguage}

%----------------------------------------------------------------------------------------
%	AGRADECIMENTOS (opcional)
%----------------------------------------------------------------------------------------

\begin{acknowledgements}
% TODO (opcional): agradecimentos.
\end{acknowledgements}

%----------------------------------------------------------------------------------------
%	ÍNDICES
%----------------------------------------------------------------------------------------

\tableofcontents

\listoffigures

\listoftables

\iflanguage{portuguese}{
\renewcommand{\listalgorithmname}{Lista de Algor\'itmos}
}
\listofalgorithms
\addchaptertocentry{\listalgorithmname}

%----------------------------------------------------------------------------------------
%	SÍMBOLOS
%----------------------------------------------------------------------------------------

\begin{symbols}{lll} % Lista de símbolos (três colunas)

$r_t$ & retorno logarítmico de um ativo no dia $t$ & --- \\
$\mu_t$ & média móvel dos retornos & --- \\
$\sigma_t$ & desvio-padrão móvel dos retornos & --- \\
$z_t$ & z-score do retorno no dia $t$ & --- \\
$k$ & limiar de anomalia (em desvios-padrão) & --- \\
$w$ & comprimento da janela móvel & \si{\day} \\

\end{symbols}

%----------------------------------------------------------------------------------------
%	ACRÓNIMOS
%----------------------------------------------------------------------------------------

\newcommand{\listacronymname}{List of Acronyms}
\iflanguage{portuguese}{
\renewcommand{\listacronymname}{Lista de Acr\'onimos}
}

\glsresetall

\printnoidxglossary[title=\listacronymname,type=\acronymtype,style=long]

%----------------------------------------------------------------------------------------
%	FIM DO FRONT MATTER
%----------------------------------------------------------------------------------------

\mainmatter
\pagestyle{thesis}
